\documentclass[../../main.tex]{subfiles}
\graphicspath{{../../image/}{../../figure/}}
\begin{document}

\section{正则化}\label{section:ch07-正则化}

在神经网络前向传播过程中，对计算得到的变量进行归一化处理可以避免网络处理过大或过小的值。尽管原则上神经网络的权重和偏置可以适应各种不同输入和隐藏变量的值，但在实践中，归一化对于确保有效的训练至关重要。根据归一化操作所针对的对象不同，我们将归一化技术分为三种类型：数据归一化、批量归一化和层归一化。

\subsection{数据归一化}\label{section:ch07-数据归一化}

有时会遇到这样的数据集，其中不同输入变量的取值范围差异很大。例如在医疗健康数据中，患者身高可能以米为单位测量（如~1.8~米），而他们的血小板计数可能以每微升血液中的血小板数量来衡量（如每微升~300\,000~个血小板）。这种差异会让梯度下降训练过程更困难：考虑一个具有两个权重的单层回归网络，其两个对应输入变量具有截然不同取值范围。其中一个权重的微小变化，相比于另一个权重的类似变化，会对输出乃至误差函数产生大得多的影响。这会产生一个沿着不同坐标轴曲率差异极大的误差曲面，如图~7.3~所示。

因此对于连续型输入变量，将输入值重新缩放使其具有相似取值范围通常是非常有益的。这很容易做到：首先分别计算每个输入变量的均值和方差
\begin{align}
\mu_i = \frac{1}{N}\sum_{n=1}^{N} x_{ni},
\label{align:ch07-input-mean}
\end{align}
\begin{align}
\sigma_i^2 = \frac{1}{N}\sum_{n=1}^{N}(x_{ni}-\mu_i)^2,
\label{align:ch07-input-var}
\end{align}
这是在开始任何训练之前都需要进行的计算。然后使用以下方法缩放输入值：
\begin{align}
\tilde{x}_{ni} = \frac{x_{ni}-\mu_i}{\sigma_i},
\label{align:ch07-norm-input}
\end{align}
这样缩放后的值 $\tilde{x}_{ni}$ 具有零均值和单位方差（见习题~7.14）。请注意，必须使用相同的 $\mu_i$ 值和 $\sigma_i$ 值来预处理任何开发集、验证集或测试集数据，以确保对所有输入都以相同方式进行缩放。

\subsection{批量归一化}\label{section:ch07-批量归一化}

我们已经了解了数据归一化的重要性，类似思路也可以应用于深度网络中每一隐藏层的变量。如果某个特定隐藏层的激活值范围存在很大差异，那么将这些值归一化为零均值和单位方差应该可以简化后续网络层的学习问题。然而与数据归一化可在训练开始前一次性完成不同，隐藏单元值的归一化需要在训练过程中每次权重更新后都重复进行。这种方法称为\textbf{批量归一化}（batch normalization）~\citep{IoffeSzegedy2015}。

引入批量归一化也是为了解决\textbf{梯度消失}和\textbf{梯度爆炸}问题，它们通常出现在训练非常深的神经网络时。根据微积分中的链式法则，误差函数 $E$ 相对于网络第一层中参数的梯度由下式给出：
\begin{align}
\frac{\partial E}{\partial w_{ij}^{(1)}} = \sum_{k}\cdots\sum_{j'} \frac{\partial E}{\partial z^{(K)}_{k}}\cdots\frac{\partial z^{(1)}_{j}}{\partial a^{(1)}_{i}}\frac{\partial a^{(1)}_{i}}{\partial w^{(1)}_{ij}},
\label{align:ch07-chain-deep}
\end{align}
其中 $z^{(K)}_{k}$ 表示第 $K$ 层中节点 $k$ 的激活值，式~\eqref{align:ch07-chain-deep}~右侧每个偏导数（参见~8.1.5~小节）代表那一层雅可比矩阵的元素。如果这些偏导数项中大多数的绝对值小于~1，那么大量此类项的乘积将趋于 $0$；反之如果大多数项的绝对值大于~1，乘积将趋于无穷大。因此随着网络深度增加，误差函数梯度可能变得非常大或非常小。批量归一化方法在很大程度上解决了这个问题。

为明确批量归一化的具体定义，考虑多层网络中某个特定层。该层中每个隐藏单元都对其输入的预激活值 $z_j = h(a_j)$ 计算一个非线性函数，因此我们可以选择对预激活值 $a_j$ 或激活值 $z_j$ 进行归一化。在实践中二者均可，这里以归一化预激活值 $a_j$ 为例说明过程。由于权重值在每个小批量之后都会更新，因此需要对每个小批量分别归一化。具体来说，对于一个大小为 $K$ 的小批量，定义
\begin{align}
\mu_B = \frac{1}{K}\sum_{n=1}^{K} a_n,
\label{align:ch07-batch-mean}
\end{align}
\begin{align}
\sigma_B^2 = \frac{1}{K}\sum_{n=1}^{K}(a_n-\mu_B)^2,
\label{align:ch07-batch-var}
\end{align}
\begin{align}
\hat{a}_n = \frac{a_n-\mu_B}{\sqrt{\sigma_B^2+\delta}},
\label{align:ch07-batch-norm}
\end{align}
其中对 $n=1,\cdots,K$ 的求和操作针对小批量中的元素进行；$\delta$ 是一个很小常数，引入它是为了避免当 $\sigma_B^2$ 很小时出现因为除数趋于零而导致的各种数值问题。

通过对网络中某一层的预激活值进行归一化，我们减少了该层参数自由度，从而降低了它的表达能力。我们可以通过将预激活值重新缩放为具有平均值 $\beta_j$ 和标准差 $\gamma_j$ 来补偿这一点：
\begin{align}
y_j = \gamma_j\,\hat{a}_n + \beta_j,
\label{align:ch07-batch-scale}
\end{align}
其中 $\beta_j$ 和 $\gamma_j$ 是可学习的自适应参数，它们与网络权重和偏置一起通过梯度下降法进行联合学习。这些可学习参数是批量归一化与输入数据归一化的一个关键区别：在式~\eqref{align:ch07-batch-scale}~中，它们直接由独立参数 $\beta_j$ 和 $\gamma_j$ 决定，这些参数在梯度下降过程中更容易学习。

式~\eqref{align:ch07-batch-scale}~可能看起来只是抵消了批量归一化的效果，因为均值和方差现在又可以适应任意值了。然而关键区别在于参数在训练过程中的演变方式：对于原始网络，一个小批量的均值和方差由该层所有权重和偏置的复杂函数决定；而在式~\eqref{align:ch07-batch-scale}~中，它们直接由独立参数 $\beta_j$ 和 $\gamma_j$ 决定。式~\eqref{align:ch07-batch-norm}~和式~\eqref{align:ch07-batch-scale}~描述了一个相对于小批量数据点的、由可学习参数 $\beta$ 和 $\gamma$ 决定的可微变量转换。这可视为神经网络中的附加层，因此每个标准隐藏层后面都可以加一个批量归一化层。

当完成网络训练、想要对新数据进行预测时，我们并没有训练时的小批量数据可用，也不能仅从一个数据样本确定均值和方差。为解决这个问题，原则上可以在对权重和偏置进行最后更新之后，针对每个层在整个训练集上计算 $\mu$ 和 $\sigma^2$；然而为计算这些统计量就需要对整个数据集再处理一遍，计算成本太高。于是选择在整个训练阶段计算移动平均值：
\begin{align}
\mu^{(\tau)} = \alpha\,\mu^{(\tau-1)} + (1-\alpha)\mu_B,
\label{align:ch07-runmean}
\end{align}
\begin{align}
\sigma^{2(\tau)} = \alpha\,\sigma^{2(\tau-1)} + (1-\alpha)\sigma_B^2,
\label{align:ch07-runvar}
\end{align}
其中 $0\le\alpha\le 1$。这些移动平均值在训练期间不起作用，仅用于在推理阶段处理新数据点。

尽管批量归一化在实践中非常有效，但其内在原理目前还不明确。批量归一化提出的原因是：人们观察到网络中较早层权重的更新会改变后续层所看到的数值分布，这种现象称为\textbf{内部协变量偏移}（internal covariate shift）。然而后来的研究~\citep{Santurkar2018}表明，协变量偏移并不是一个关键因素，批量归一化所带来训练改进实际上在于它平滑了误差函数的崎岖程度。

\subsection{层归一化}\label{section:ch07-层归一化}

在批量归一化中，如果批量规模太小，那么均值和方差的估计会包含太多噪声。另外，对于非常大的训练集，小批量可能会被分散在不同 GPU 上处理，使得跨小批量的全局归一化效率低下。一种替代方案是分别针对每个数据点对隐藏单元的值进行归一化，而不是在小批量中的样本间进行归一化。这称为\textbf{层归一化}（layer normalization）~\citep{BaKirosHinton2016}。

层归一化对每个样本在同一层的所有隐藏单元之间归一化，从而不依赖小批量统计，对批量大小不敏感，在循环神经网络和小批量训练场景中尤为有效。

\end{document}
